\section{Titre II --- La fin du contrat}

\article[absolue]{Causes de fin}\label{art:causes}
Le contrat ne prend fin que par une rupture pendant la période d'essai, une
résiliation ordinaire, une résiliation immédiate pour juste motif, un accord
écrit des parties ou le décès du salarié. L'écoulement du temps n'y met
jamais fin.
\lean{CauseFin}
\lean{execute}
\theoreme{fin_justifiee}
\theoreme{pas_de_terme}

\article{Délais et ancienneté}\label{art:delais}
I. --- Les préavis se comptent en jours. Une notification prend effet le jour
de sa réception ; le préavis court à partir du lendemain. Le dernier jour du
préavis est le dernier jour du contrat.

II. --- L'ancienneté se compte en jours depuis le début du contrat. Une année
révolue compte trois cent soixante-cinq jours ; le nombre de mois révolus est
la partie entière de douze fois le nombre de jours, divisé par trois cent
soixante-cinq.
\lean{anneesRevolues}
\lean{moisRevolus}

\article[relative]{Rupture pendant l'essai}\label{art:rupture-essai}
Pendant la période d'essai, chaque partie peut mettre fin au contrat à tout
moment, sans motif, moyennant le préavis prévu à
l'article~\ref{art:resiliation}. Le préavis court même s'il se termine après
la fin de la période d'essai.
\lean{dispPreavisEmployeur}
\lean{dispPreavisSalarie}
\theoreme{fin_ruptureEssai}

\article[relative]{Résiliation ordinaire}\label{art:resiliation}
I. --- Après la période d'essai, chaque partie peut résilier le contrat à tout
moment, sans motif, moyennant un préavis. Aucune résiliation ordinaire ne peut
être notifiée pendant la période d'essai.

II. --- La durée du préavis, en jours, est la partie entière de la plus petite
des quatre quantités suivantes, où $j$ est l'ancienneté du salarié en jours au
jour de la réception de la notification :
\[
  7 + \frac{23}{30}\,j, \qquad
  30 + \frac{6}{67}\,(j - 30), \qquad
  60 + \frac{3}{292}\,(j - 365), \qquad
  90.
\]
Le préavis est ainsi de sept jours au début du contrat, de trente jours après
trente jours, de soixante jours après un an et de quatre-vingt-dix jours après
neuf ans ; entre ces dates, il croît régulièrement, d'au plus un jour par jour
d'ancienneté.

III. --- Un accord d'entreprise ou le contrat peut allonger le préavis dû par
l'employeur et raccourcir celui dû par le salarié
(article~\ref{art:relatives}).
\lean{preavis}
\lean{dispPreavisEmployeur}
\lean{dispPreavisSalarie}
\theoreme{preavis_monotone}
\theoreme{fin_preavis_succ}
\theoreme{fin_resiliationOrdinaire}
\theoreme{duree_minimale}
\theoreme{preavis_ordinaire_apres_essai}

\article[relative]{Périodes de protection}\label{art:protection}
I. --- Après la période d'essai, pendant une incapacité de travail du salarié
ou son congé de maternité :
\begin{itemize}
  \item la résiliation ordinaire notifiée par l'employeur est nulle ;
  \item le préavis d'une résiliation ordinaire de l'employeur déjà notifiée
        est suspendu, et reprend à la fin de la période de protection.
\end{itemize}

II. --- La protection ne fait obstacle ni à la démission du salarié, ni à
une résiliation immédiate pour juste motif.
\lean{suspendu}
\theoreme{protection_nullite}
\theoreme{protection_suspend}
\theoreme{protection_salarie_libre}

% À rédiger (phase 6) : congé abusif.

\article[absolue]{Résiliation immédiate pour juste motif}\label{art:resiliation-immediate}
Chaque partie peut résilier le contrat sans préavis pour juste motif. Le
contrat prend fin le jour de la réception de la résiliation. Le juste motif
est une notion ouverte, appréciée par le juge.
\lean{Interpretation.justeMotif}
\theoreme{fin_resiliationImmediate}

\article[absolue]{Accord écrit et décès}\label{art:accord-deces}
I. --- Les parties peuvent mettre fin au contrat d'un commun accord, par écrit,
sans préavis.

II. --- Le contrat prend fin de plein droit au décès du salarié.
\lean{CauseFin.accordCommun}
\lean{CauseFin.deces}

\article[relative]{Motivation sur demande}\label{art:motivation}
La partie qui résilie le contrat en indique le motif par écrit si l'autre
partie le demande. L'absence de motivation n'affecte pas la validité de la
résiliation.
\horsformalisation{obligation d'information, sans effet sur la machine à états}

\article{Salaire de référence}\label{art:salaire-reference}
I. --- Le salaire mensuel est la rémunération brute totale versée au titre
d'un mois, primes et heures supplémentaires comprises. Le taux d'activité d'un
mois est la fraction de ce mois pendant laquelle le salarié a travaillé ou a été
rémunéré, congés payés compris ; il est compris entre zéro et un.

II. --- Tous les mois du contrat sont pris en compte, du plus ancien au plus
récent. La masse des salaires et la masse d'activité sont nulles avant le
premier mois. À chaque mois, la masse des salaires devient la somme d'un
vingt-quatrième du salaire de ce mois et des vingt-trois vingt-quatrièmes de la
masse des salaires précédente ; la masse d'activité devient la somme d'un
vingt-quatrième du taux d'activité de ce mois et des vingt-trois
vingt-quatrièmes de la masse d'activité précédente.

III. --- Le salaire de référence est le quotient de la masse des salaires par la
masse d'activité. Il est nul si la masse d'activité est nulle.
\lean{MoisTravaille}
\lean{salaireReference}
\theoreme{salaireReference_encadrement}
\theoreme{salaireReference_sensibilite}
\theoreme{salaireReference_prorata}
\theoreme{salaireReference_mois_proche}

\article[relative]{Indemnité de rupture}\label{art:indemnite}
I. --- En cas de résiliation ordinaire du contrat par l'employeur, le
salarié a droit à une indemnité de rupture.

II. --- Le montant de l'indemnité est
\[
  I = \left\lfloor SR \times 6 \times S\!\left(\frac{\min(m,\,240)}{240}\right) \right\rfloor
  \qquad\text{avec}\qquad S(t) = 3t^2 - 2t^3,
\]
où $SR$ est le salaire de référence défini à l'article~\ref{art:salaire-reference}
et $m$ l'ancienneté du salarié en mois révolus au dernier jour du contrat
(article~\ref{art:delais}).
Le montant est arrondi au centime d'euro inférieur. Aucun calcul intermédiaire
n'est arrondi.

III. --- L'indemnité n'est due ni en cas de démission, ni en cas de rupture
pendant la période d'essai, ni en cas de résiliation immédiate pour juste motif,
ni en cas d'accord écrit des parties ou de décès du salarié.
\lean{indemniteRupture}
\theoreme{indemniteRupture_le_plafond}
\theoreme{indemnite_applicable_le_plafond}
\theoreme{indemniteRupture_manipulation}
\theoreme{indemniteRupture_eq_zero}
\theoreme{indemnite_mono_anciennete}
\theoreme{indemnite_saut_mensuel}
\theoreme{indemnite_ge_legal}

\begin{center}
\footnotesize
\begin{tabular}{lccccc}
  \multicolumn{6}{c}{\textit{Valeurs indicatives, en mois de salaire de référence}} \\[0.3em]
  \hline
  Ancienneté & 1 an & 5 ans & 10 ans & 15 ans & 20 ans et plus \\
  Indemnité & 0,04 & 0,94 & 3 & 5,06 & 6 \\
  \hline
\end{tabular}
\end{center}
